% GENERATED FILE. Edit canonical lessons, then run npm run generate:books.
% canonical-source-hash: fnv1a64:fce75e73f84bd790
% canonical-lessons: AR-C426-jamarik, AR-C426-fattasha, AR-C426-khatama, AR-C426-afragha, AR-C426-sadara

\chapter{Through customs}
\label{ch:ar-through-customs}
\hlchaptermodality{426}

\begin{chapteropening}
\textbf{By the end of this chapter:} \emph{I can say customs, to search, to stamp, to empty out and to confiscate.}

Complete the last lesson of chapter 426: I can say customs, to search, to stamp, to empty out and to confiscate.
\end{chapteropening}

\section[\texorpdfstring{\ar{جمارك} (jamārik)}{جمارك (jamārik)}]{\ar{جمارك} (jamārik) — customs (at a border)}

\label{lesson:AR-C426-jamarik}



\begin{quote}
Before the new one: say the Arabic for to tighten, then the Arabic for arrival.
\end{quote}

\subsection*{You'll want to know: \ar{جمارك}}
\textbf{\ar{جمارك}} — \emph{jamārik} — \textquotedblleft{}customs (at a border)\textquotedblright{}.

\begin{quote}
\ar{وقفنا} \ar{في} \ar{الجمارك} \ar{مدة} \ar{طويلة}.
\end{quote}

\begin{quote}
\emph{waqafnā fī al-jamārik mudda ṭawīla.}
\end{quote}

\begin{quote}
We stood in customs for a long time.
\end{quote}

The singular \ar{جمرك} (\emph{jumruk}) names one customs post; the plural names the service.

\begin{cousinweb}
\textbf{\ar{جمارك}} (\emph{jamārik}) is usually traced to Turkish \emph{gümrük}, which goes back to Greek \emph{kommerkion}, a duty on trade.
\end{cousinweb}

\begin{grammarlens}[title={What you've built}]
One more word for getting around.
\end{grammarlens}

\subsection*{Your turn}
\begin{itemize}
\raggedright
  \item \emph{Say it:} \emph{jamārik}
  \item \emph{Say it:} \emph{jamārik}, once more
  \item \emph{Say it:} \emph{waqafnā fī al-jamārik mudda ṭawīla.} — We stood in customs for a long time.
  \item \emph{Recall:} say the Arabic for to tighten, then the Arabic for arrival, then say \emph{jamārik} again
\end{itemize}

\begin{tcolorbox}[breakable,colback=teal!4,colframe=teal!35!black,title={Before you move on}]
What does \ar{جمارك} mean? (\textquotedblleft{}Customs (at a border)\textquotedblright{}.) Say it once more.
\end{tcolorbox}
\section[\texorpdfstring{\ar{فتش} (fattasha)}{فتش (fattasha)}]{\ar{فتش} (fattasha) — to search; to inspect}

\label{lesson:AR-C426-fattasha}



\begin{quote}
Before the new one: say the Arabic for arrival, then the Arabic for customs.
\end{quote}

\subsection*{You'll want to know: \ar{فتش}}
\textbf{\ar{فتش}} — \emph{fattasha} — \textquotedblleft{}to search; to inspect\textquotedblright{}.

\begin{quote}
\ar{فتش} \ar{الموظف} \ar{حقيبتي} \ar{في} \ar{الجمارك}.
\end{quote}

\begin{quote}
\emph{fattasha al-muwaẓẓaf ḥaqībatī fī al-jamārik.}
\end{quote}

\begin{quote}
The officer searched my bag at customs.
\end{quote}

Present: \ar{يفتش}, \emph{yufattish}.

\begin{cousinweb}
\textbf{\ar{فتش}} (\emph{fattasha}) is built on the native Arabic root \emph{f-t-sh}, \textquotedblleft{}to search\textquotedblright{} (Lane). It is Form II, \textquotedblleft{}to search thoroughly\textquotedblright{}.
\end{cousinweb}

\begin{grammarlens}[title={What you've built}]
One more word for getting around.
\end{grammarlens}

\subsection*{Your turn}
\begin{itemize}
\raggedright
  \item \emph{Say it:} \emph{fattasha}
  \item \emph{Say it:} \emph{fattasha}, once more
  \item \emph{Say it:} \emph{fattasha al-muwaẓẓaf ḥaqībatī fī al-jamārik.} — The officer searched my bag at customs.
  \item \emph{Recall:} say the Arabic for arrival, then the Arabic for customs, then say \emph{fattasha} again
\end{itemize}

\begin{tcolorbox}[breakable,colback=teal!4,colframe=teal!35!black,title={Before you move on}]
What does \ar{فتش} mean? (\textquotedblleft{}To search; to inspect\textquotedblright{}.) Say it once more.
\end{tcolorbox}
\section[\texorpdfstring{\ar{ختم} (khatama)}{ختم (khatama)}]{\ar{ختم} (khatama) — to stamp; to seal}

\label{lesson:AR-C426-khatama}



\begin{quote}
Before the new one: say the Arabic for customs, then the Arabic for to search.
\end{quote}

\subsection*{You'll want to know: \ar{ختم}}
\textbf{\ar{ختم}} — \emph{khatama} — \textquotedblleft{}to stamp; to seal\textquotedblright{}.

\begin{quote}
\ar{ختم} \ar{الموظف} \ar{جوازي} \ar{عند} \ar{الوصول}.
\end{quote}

\begin{quote}
\emph{khatama al-muwaẓẓaf jawāzī ʿinda al-wuṣūl.}
\end{quote}

\begin{quote}
The officer stamped my passport on arrival.
\end{quote}

Present: \ar{يختم}, \emph{yakhtim}.

The noun \emph{khatm} is a stamp; \emph{khatama} also means to finish.

\begin{cousinweb}
\textbf{\ar{ختم}} (\emph{khatama}) is built on the native Arabic root \emph{kh-t-m}, \textquotedblleft{}to seal\textquotedblright{} (Lane). It is the root of \emph{khātam}, \textquotedblleft{}a ring\textquotedblright{}, once a seal ring.
\end{cousinweb}

\begin{grammarlens}[title={What you've built}]
One more word for getting around.
\end{grammarlens}

\subsection*{Your turn}
\begin{itemize}
\raggedright
  \item \emph{Say it:} \emph{khatama}
  \item \emph{Say it:} \emph{khatama}, once more
  \item \emph{Say it:} \emph{khatama al-muwaẓẓaf jawāzī ʿinda al-wuṣūl.} — The officer stamped my passport on arrival.
  \item \emph{Recall:} say the Arabic for customs, then the Arabic for to search, then say \emph{khatama} again
\end{itemize}

\begin{tcolorbox}[breakable,colback=teal!4,colframe=teal!35!black,title={Before you move on}]
What does \ar{ختم} mean? (\textquotedblleft{}To stamp; to seal\textquotedblright{}.) Say it once more.
\end{tcolorbox}
\section[\texorpdfstring{\ar{أفرغ} (afragha)}{أفرغ (afragha)}]{\ar{أفرغ} (afragha) — to empty out; to unpack}

\label{lesson:AR-C426-afragha}



\begin{quote}
Before the new one: say the Arabic for to search, then the Arabic for to stamp.
\end{quote}

\subsection*{You'll want to know: \ar{أفرغ}}
\textbf{\ar{أفرغ}} — \emph{afragha} — \textquotedblleft{}to empty out; to unpack\textquotedblright{}.

\begin{quote}
\ar{أفرغ} \ar{المسافر} \ar{حقيبته} \ar{أمام} \ar{الموظف}.
\end{quote}

\begin{quote}
\emph{afragha al-musāfir ḥaqībatahu amāma al-muwaẓẓaf.}
\end{quote}

\begin{quote}
The traveller emptied out his bag in front of the officer.
\end{quote}

Present: \ar{يفرغ}, \emph{yufrigh}.

\begin{cousinweb}
\textbf{\ar{أفرغ}} (\emph{afragha}) is built on the native Arabic root \emph{f-r-gh}, \textquotedblleft{}to be empty, to be free\textquotedblright{} (Lane). It is Form IV, \textquotedblleft{}to make empty\textquotedblright{}, beside \emph{fārigh}, \textquotedblleft{}empty\textquotedblright{}.
\end{cousinweb}

\begin{grammarlens}[title={What you've built}]
One more word for getting around.
\end{grammarlens}

\subsection*{Your turn}
\begin{itemize}
\raggedright
  \item \emph{Say it:} \emph{afragha}
  \item \emph{Say it:} \emph{afragha}, once more
  \item \emph{Say it:} \emph{afragha al-musāfir ḥaqībatahu amāma al-muwaẓẓaf.} — The traveller emptied out his bag in front of the officer.
  \item \emph{Recall:} say the Arabic for to search, then the Arabic for to stamp, then say \emph{afragha} again
\end{itemize}

\begin{tcolorbox}[breakable,colback=teal!4,colframe=teal!35!black,title={Before you move on}]
What does \ar{أفرغ} mean? (\textquotedblleft{}To empty out; to unpack\textquotedblright{}.) Say it once more.
\end{tcolorbox}
\section[\texorpdfstring{\ar{صادر} (ṣādara)}{صادر (ṣādara)}]{\ar{صادر} (ṣādara) — to confiscate}

\label{lesson:AR-C426-sadara}



\begin{quote}
Before the new one: say the Arabic for to stamp, then the Arabic for to empty out.
\end{quote}

\subsection*{You'll want to know: \ar{صادر}}
\textbf{\ar{صادر}} — \emph{ṣādara} — \textquotedblleft{}to confiscate\textquotedblright{}.

\begin{quote}
\ar{صادر} \ar{موظف} \ar{الجمارك} \ar{العسل} \ar{من} \ar{حقيبتي}.
\end{quote}

\begin{quote}
\emph{ṣādara muwaẓẓaf al-jamārik al-ʿasal min ḥaqībatī.}
\end{quote}

\begin{quote}
The customs officer confiscated the honey from my bag.
\end{quote}

Present: \ar{يصادر}, \emph{yuṣādir}.

\begin{cousinweb}
\textbf{\ar{صادر}} (\emph{ṣādara}) is built on the native Arabic root \emph{ṣ-d-r}, \textquotedblleft{}to come forth\textquotedblright{} (Lane). It is Form III.
\end{cousinweb}

\begin{grammarlens}[title={What you've built}]
One more word for getting around.
\end{grammarlens}

\subsection*{Your turn}
\begin{itemize}
\raggedright
  \item \emph{Say it:} \emph{ṣādara}
  \item \emph{Say it:} \emph{ṣādara}, once more
  \item \emph{Say it:} \emph{ṣādara muwaẓẓaf al-jamārik al-ʿasal min ḥaqībatī.} — The customs officer confiscated the honey from my bag.
  \item \emph{Recall:} say the Arabic for to stamp, then the Arabic for to empty out, then say \emph{ṣādara} again
\end{itemize}

\begin{tcolorbox}[breakable,colback=teal!4,colframe=teal!35!black,title={Before you move on}]
What does \ar{صادر} mean? (\textquotedblleft{}To confiscate\textquotedblright{}.) Say it once more.
\end{tcolorbox}
